% Audited research proof; stronger bounds appear in the following refinement.
% Include after random-host-weighted-matching.tex and random-hosts-phase4.tex.
\section{Using the wheel rims as additional short cycles}
\label{sec:rhost-seventy-nine}

The previous wheel construction did not use its rim cycle to fill a distinct required length.
We now use a small proportion of longer rims for this purpose.
The resulting coefficient is $79/72<10/9$.

\begin{lemma}[A weighted wheel packing]\label{lem:rhost-wheel-profit}
Let $h\to\infty$ through multiples of nine.
There is a list of $O(h)$ traceable wheel subdivisions with the following properties.
Their non-rim cycle lengths form pairwise disjoint twelve-element subsets of $[h+1,2h]$.
Some wheels have a marked rim length in $[3,h/3+1]$.
The marked rim lengths are pairwise distinct.
If $a$ wheels have no marked rim and $b$ wheels have a marked rim, then
\[
 8a+9b\ge\left(\frac{49}{72}-o(1)\right)h.
\]
Every wheel has $O(h)$ edges, shortest cycle length at least $2\lceil\sqrt h\rceil$,
and at least $(r-1)h$ edges in every connected core of rank $r\in\{2,3,4\}$.
\end{lemma}

\begin{proof}
Put $s=\lceil\sqrt h\rceil$.
Use the spoke distribution from Lemma~\ref{lem:rhost-wheel-lengths}:
$U$ is uniform on $\{s,\ldots,2s\}$, the four deficits have sum $h+U$ and distribution
\eqref{eq:rhost-wheel-distribution}, and spoke $i$ has length $h+s-X_i$.
Independently choose one of two types.

With probability $2/3$, use the previous small rim $(1,L,1,M)$, where $L,M$ are independent and uniform on $\{s,\ldots,2s\}$.
Ignore its rim length and give the resulting twelve-element length set value eight.
With probability $1/3$, choose
\[
 R\ \text{uniform on }\{6s,\ldots,\lfloor h/36\rfloor\}.
\]
Choose $L$ uniformly from the integer interval with center $\lfloor(R-2)/2\rfloor$ and radius $\lfloor s/2\rfloor$.
Put $M=R-2-L$ and use the rim $(1,L,1,M)$.
Mark its rim length $R$ and give the thirteen-element length set value nine.
For sufficiently large $h$, these rim arms are positive.
The marked rim length is in $[3,h/3+1]$.

Discard outcomes with a nonpositive spoke, repeated non-rim cycle lengths, or a non-rim cycle outside $[h+1,2h]$.
For the small-rim type, the discarded probability is $o(1)$.
We estimate it for the large-rim type.
The first two causes still have probability $o(1)$.
The spoke calculation is unchanged.
For repeats, the codegree calculation below applies.

For a fixed spoke pair with rim-arc lengths $a,a'$, a cycle exceeds $2h$ only if
\[
 X_i+X_j<2s+\max\{a,a'\}.
\]
The pair deficit is uniform on $\{0,\ldots,h+U\}$.
The two pairs $01,23$ have largest rim arc $R-1$.
The other four pairs have largest rim arc at most $R/2+s/2+2$.
A union bound over these six pairs gives
\[
 \Pr(\text{some non-rim cycle exceeds }2h\mid R)
 \le\frac{4R+14s+6}{h}.
\]
Every non-rim cycle has length at least $h+1$.
Since $\mathbb E R=h/72+O(s)$, the retained probability for the large-rim type is at least $17/18-o(1)$.
Thus the expected retained value is at least
\begin{equation}
 8\cdot\frac23+9\cdot\frac13\left(\frac{17}{18}-o(1)\right)
 =\frac{49}{6}-o(1).
 \label{eq:rhost-wheel-value}
\end{equation}

For every upper-interval length, its weighted degree is at most $12/h$.
This follows from the uniform pair-deficit marginal, separately for each of the twelve labels.
For a marked rim length, the weighted degree is at most
\[
 \frac{1/3}{\lfloor h/36\rfloor-6s+1}
 =\frac{12+o(1)}h.
\]
All weighted codegrees are $O(h^{-3/2})$.
For two upper-interval lengths, the shared-index argument from Lemma~\ref{lem:rhost-wheel-lengths} is unchanged.
For complementary spoke pairs, the independent variable $U$ gives the same bound.
For two arcs of one spoke pair, their difference is smoothed either by $R$, or by the interval of $\Theta(s)$ choices for $L$.
For a rim length and an upper-interval length, first fixing $R$ costs $O(1/h)$.
Fixing the pair deficit then costs at most $1/h$.
Two distinct marked rim lengths never occur in the same outcome.
These bounds also show that repeated upper-interval labels have probability $o(1)$.

Apply Lemma~\ref{lem:rhost-weighted-matching} to these weighted sets, with values eight and nine.
The maximum weighted degree is $(12+o(1))/h$ and the weighted total value is at least $49/6-o(1)$.
It gives a disjoint collection with total value at least $(49/72-o(1))h$.
Choose one retained wheel metric for each selected set.

All chosen metrics satisfy the claimed density and traceability properties.
Indeed, their total spoke length is $3h+4s-U$, every spoke is at most $h+s$,
and every two spokes have total length at least $h$.
Both rim types have length at least $2s$ and at most $O(h)$.
The core classification and spanning-path deletion from Lemma~\ref{lem:rhost-wheel-lengths} apply verbatim.
\end{proof}

\begin{theorem}[An upper coefficient of $79/72$]\label{thm:rhost-seventy-nine}
For the same $p_*$ as in Theorem~\ref{thm:rhost-scale}, and every $p\ge p_*$,
with high probability
\[
 \operatorname{Pex}(G(n,p))
 \le\log_2n+\left(\frac{79}{72}+o(1)\right)\frac{\ln n}{\ln\ln n}.
\]
\end{theorem}

\begin{proof}
Use the exposure scheme of Theorem~\ref{thm:rhost-ten-ninths}, with $h=9\lceil d/9\rceil$.
Start with single-cycle blocks of lengths $3,\ldots,h/3+1$,
the theta partition for $[h/3+2,h-2]$, and single-cycle blocks of lengths $h-1,\ldots,2h$.
Their total rank is $16h/9+O(1)$.

Take the wheel list from Lemma~\ref{lem:rhost-wheel-profit}.
For each wheel, replace its twelve single-cycle blocks in the upper interval by that wheel.
If its rim is marked, also remove the single-cycle block of that rim length.
The disjointness of the chosen length sets makes all these replacements compatible.
A wheel has rank four, so an unmarked wheel saves eight and a marked wheel saves nine.
The resulting short-block rank is at most
\[
 \frac{16h}9-\left(\frac{49}{72}-o(1)\right)h
 =\left(\frac{79}{72}+o(1)\right)h.
\]
Add $o(d)$ single-cycle blocks for the remaining required lengths through $T=2(d+b+1)$.

The wheel overlap proof in Lemma~\ref{lem:rhost-wheel-embeddings} uses only the stated core thresholds,
the lower bound on the shortest cycle, and the $O(h)$ total edge bound.
It therefore applies to both rim types in the present list.
The theta embedding, connector chain, remaining exposures, and final excess count are unchanged.
Since $h=d+O(1)$ and $b,m=o(d)$, the asserted bound follows.
\end{proof}
